\chapter{尾部与压力风险提示}
\begin{enumerate}[leftmargin=*]
  \item 先判断问题是概率尾部还是条件冲击。
  \item 把尾部积分拆成完整观察与边界概率质量。
  \item 令冲击为 $sx$，把二阶损失写成 $s$ 的二次式。
  \item 直接计算 $n(1-\alpha)$，不要只看总样本量。
  \item 保持同一风险因子和估值时点，再比较重估层级。
  \item 99\% 的四点样本只有 0.04 个有效尾部观察。
  \item 它是条件阈值，不携带冲击方向的概率模型。
  \item 报告必须能从冻结 fixture、oracle 和命令重建。
\end{enumerate}
